% kernel/engineering_identity.tex

\chapter{Identity Through Transition Change}
\label{chap:engineering-identity}

The proposed transition now refers to one alignment. The engineer changes the
length of the entry half-wave and recomputes the coupled parameters. Geometry
and derived pose change. The subject need not.

\section{Object continuity and state succession}

\begin{definition}[Engineering object identity]
\label{def:engineering-object-identity}
Engineering Object Identity is the basis by which an Engineering Object
remains distinguishable throughout its engineering lifecycle.
\end{definition}

Object identity answers \emph{which enduring alignment?}; state identity
answers \emph{which qualified condition or revision of it?} The accepted
predecessor and the calculated proposal may therefore refer to the same
alignment while remaining different states. A recorded derivation relation
connects them without overwriting the earlier evidence.

A parameter edit normally supplies a candidate successor state: the alignment
reference, longitudinal order, and traceable predecessor remain present while
the parameterized construction changes. This is an explanatory case, not a
universal lifecycle rule. Acceptance of the successor still requires the
applicable authority.

\section{Constructive identity makes the edit intelligible}

The Berlinish expression is not merely display geometry. Its ordered
components, intrinsic longitudinal parameterization, function-family choices,
and boundary relations make the construction inspectable. Replacing the
clothoid core by another admissible family is therefore a constructive change,
even if a sampled polyline happens to remain close to the old one. Resampling
or transforming that polyline is only a representational change unless the
engineering construction is also changed.

\begin{definition}[Identity composition]
\label{def:identity-composition}
Identity composition is the coherent combination of stable identity aspects
and canonical relations required to distinguish an object without reducing
identity to a representation-dependent attribute.
\end{definition}

For an alignment, Constructive Identity contributes the intrinsic sequence and
longitudinal construction by which the elements form the alignment. Alignment
Identity draws on that constructive distinction; it is not a CRS identifier,
station label, database key, or generated geometry. Operational stationing may
address positions along the object, but it does not replace intrinsic
longitudinal parameterization as an identity-bearing construction.

\section{When does change create another object?}

Some cases are clear. Coordinate transformation, serialization, rendering,
and regeneration of derived geometry do not by themselves replace the
alignment. Preserving the source record while deriving an editable model also
does not erase the source identity or provenance. At the other end, an
explicitly commissioned replacement alignment can be governed as another
object even if it reuses much geometry.

Between those cases lies a genuine boundary. Replacing a transition component
may be admitted as a successor state or may participate in a replacement
object, depending on the object-type-specific sameness conditions and the
authorized engineering act. The active Kernel does not prescribe a universal
threshold for that judgment. Geometric distance, a software key, or the solver
operation cannot decide it.

\section{Provenance preserves the path through change}

The candidate records its predecessor, the calculation service and version,
input conditions, chosen function families and parameters, resulting
derivations, responsible actor, and time. If the candidate is rejected, this
evidence remains. If it is accepted, the decision relates the new accepted
state to its predecessor rather than rewriting history.

Identity now tells us which object and state the proposal concerns. It still
does not tell us where the transition is metrically meaningful or whether its
evidence is sufficient for a particular task. Those dependencies establish
applicability.
